\documentclass[12pt]{article}
\usepackage[T1]{fontenc}
\usepackage[polish]{babel}
\usepackage[utf8]{inputenc}
\newcommand{\BibTeX}{{\sc Bib}\TeX} 
\usepackage{graphicx}
\usepackage{amsfonts}
\usepackage[hidelinks]{hyperref}
\usepackage{float}

\usepackage{xcolor}
\usepackage{listings}
\lstset{
	basicstyle=\ttfamily\small,
	keywordstyle=\color{blue},
	commentstyle=\color{gray},
	stringstyle=\color{green!50!black},
	backgroundcolor=\color{black!5},
	showstringspaces=false,
	breaklines=true,
	frame=single,
	columns=flexible,
	keepspaces=true
}

\setlength{\textheight}{21cm}

\title{{\bf Zadanie nr 2 - Próbkowanie i kwantyzacja}\linebreak
	Cyfrowe Przetwarzanie Sygnałów}
\author{Bartosz Kołaciński, 251554 \and Jakub Rosiak, 251620}
\date{\today}

\begin{document}
	\clearpage\maketitle
	\thispagestyle{empty}
	\newpage
	\setcounter{page}{1}
	
	\section{Cel zadania}
	Celem zadania jest przeprowadzenie eksperymentów związanych z konwersją analogowo-cyfrową (A/C) i cyfrowo-analogową (C/A). 

	\section{Wstęp teoretyczny}

	Wymagana wiedza teoretyczna potrzebna do wykonania zadania została zaczerpnięta z dostępnej na WIKAMP instrukcji \cite{cps-zadanie2-instrukcja}.
	
	Proces konwersji sygnału analogowego do postaci cyfrowej obejmuje trzy główne etapy: próbkowanie, kwantyzacje i rekonstrukcje sygnału.
	
	W realizacji tego zadania wykorzystaliśmy próbkowanie równomierne, które polega na pobieranie próbek w równych odstępach czasu $T_s$ określonego za pomocą wzoru:
	\begin{equation}
		T_s = \frac{1}{f_s}
	\end{equation}
	
	Kwantyzacja odpowiada za odwzorowanie wartości sygnału na skończony zbiór poziomów. Rozważono dwa warianty kwantyzacji równomiernej: z obcięciem oraz z zaokrąglaniem.
	
	Rekonstrukcja sygnału cyfrowego do postaci analogowej realizowana jest trzema metodami: ekstrapolacją zerowego rzędu (ZOH), interpolacją pierwszego rzędu (FOH) oraz rekonstrukcją z wykorzystaniem funkcji sinc (SINC).
	
	W celu oceny jakości przetwarzania sygnału wykorzystano następujące miary: błąd średniokwadratowy (MSE), stosunek sygnału do szumu (SNR), szczytowy stosunek sygnału do szumu (PSNR) oraz maksymalną różnicę (MD).

\section{Eksperymenty i wyniki}

Opis wykonywanych eksperymentów. Wymagane jest ilustrowanie przeprowadzanych doświadczeń wykresami oraz tabelami.

%%%%%%%%%%%%%%%%%%%%%%%%%%%%%%%%%%%%%%%%%%%%%%%%%%%%%%%%%%%%%%%%%%%%%%%%%%%%%%%%%%%%%%%%%%%%%%%%%%%%%%%%%%%%%%%%%
% PODROZDZIAŁ PT. EKSPERYMENT NR 1 
%%%%%%%%%%%%%%%%%%%%%%%%%%%%%%%%%%%%%%%%%%%%%%%%%%%%%%%%%%%%%%%%%%%%%%%%%%%%%%%%%%%%%%%%%%%%%%%%%%%%%%%%%%%%%%%%%

\subsection{Eksperyment nr 1}

Eksperyment nr 1...\\
Identycznościowa funkcja aktywacji ma postać:
\begin{equation}
 \forall\: s \in \mathbb{R}\:\:\:\: f(s) = s
 \label{równanie funkcji identycznościowej}
\end{equation}
Jak widać z definicji (\ref{równanie funkcji identycznościowej}) funkcja ta...

\subsubsection{Założenia}

\subsubsection{Przebieg}
\newpage

\subsubsection{Rezultat}

Rezultaty badań eksperymentalnych przedstawione są w Tab. \ref{wyniki eksperymentu pierwszego}.
\begin{table}[h!]
 \caption{Rezultaty eksperymentu nr 1}
 \centering
 \vspace{0.2cm}
 \begin{tabular}{c c c c}
  \hline\hline\\[-0.4cm]
  \textbf{Przypadek} & \textbf{Metoda 1} & \textbf{Metoda 2} & \textbf{Metoda 3}\\[0.1cm]
  \hline
  \textbf{1} & 50 & 837 & 970  \\
  \textbf{2} & 47 & 877 & 230  \\
  \textbf{3} & 31 &  25 & 415  \\
  \textbf{4} & 35 & 144 & 2356 \\
  \textbf{5} & 45 & 300 & 556  \\ [0.1cm]
  \hline
 \end{tabular}
 \label{wyniki eksperymentu pierwszego}
\end{table}

\noindent Jak widać w Tab. \ref{wyniki eksperymentu pierwszego}...\newline
Graficzna interpretacja wyników z Tab. \ref{wyniki eksperymentu pierwszego} 
przedstawiona jest na wykresie Rys. \ref{rysunek do eksperymentu pierwszego} gdzie można zauważyć, że...

\noindent Jak widać z wykresu Rys. \ref{rysunek do eksperymentu pierwszego}...\newline

%%%%%%%%%%%%%%%%%%%%%%%%%%%%%%%%%%%%%%%%%%%%%%%%%%%%%%%%%%%%%%%%%%%%%%%%%%%%%%%%%%%%%%%%%%%%%%%%%%%%%%%%%%%%%%%%%
% PODROZDZIAŁ PT. EKSPERYMENT NR 2 
%%%%%%%%%%%%%%%%%%%%%%%%%%%%%%%%%%%%%%%%%%%%%%%%%%%%%%%%%%%%%%%%%%%%%%%%%%%%%%%%%%%%%%%%%%%%%%%%%%%%%%%%%%%%%%%%%

\newpage
\subsection{Eksperyment nr 2}

Eksperyment nr 2 polegał na...\\
Sigmoidalna funkcja aktywacji ma postać:
\begin{equation}
 \forall\: s \in \mathbb{R}\:\:\:\:\:\: f(s) = \frac{1}{\:\:\:1 + e^{-\beta \cdot s}\:}\:,
 \:\:\:\:\textnormal{gdzie}\:\:\beta \in \mathbb{R}_{+}
 \label{równanie funkcji sigmoidalnej}
\end{equation}
Jak widać z równania definicyjnego (\ref{równanie funkcji sigmoidalnej}) 
funkcja\footnote{ang. \textit{sigmoidal function} lub \textit{unipolar function}}
ta ma wykres przedstawiony 
na rysunku Rys. \ref{funkcja sigmoidalna}, gdzie paramater $\beta$ ...

\subsubsection{Założenia}

\subsubsection{Przebieg}

\subsubsection{Rezultat}

Rezultaty badań eksperymentalnych przedstawione są w Tab. \ref{wyniki eksperymentu drugiego}.
\vspace{-0.5cm}
\begin{table}[h!]
 \caption{Rezultaty eksperymentu nr 2}
 \centering
 \vspace{0.2cm}
 \begin{tabular}{c c c }
  \hline\hline\\[-0.4cm]
  \textbf{Przypadek} & \textbf{Metoda 1} & \textbf{Metoda 2}\\[0.1cm]
  \hline
  \textbf{1} & 50 & 837 \\
  \textbf{2} & 47 & 877 \\
  \textbf{3} & 45 & 300 \\ [0.1cm]
  \hline
 \end{tabular}
 \label{wyniki eksperymentu drugiego}
\end{table}

\noindent Jak widać w Tab. \ref{wyniki eksperymentu drugiego}...\newline
Wyniki w Tab. \ref{wyniki eksperymentu drugiego} świadczą o tym, że...

%%%%%%%%%%%%%%%%%%%%%%%%%%%%%%%%%%%%%%%%%%%%%%%%%%%%%%%%%%%%%%%%%%%%%%%%%%%%%%%%%%%%%%%%%%%%%%%%%%%%%%%%%%%%%%%%%
% PODROZDZIAŁ PT. EKSPERYMENT NR N 
%%%%%%%%%%%%%%%%%%%%%%%%%%%%%%%%%%%%%%%%%%%%%%%%%%%%%%%%%%%%%%%%%%%%%%%%%%%%%%%%%%%%%%%%%%%%%%%%%%%%%%%%%%%%%%%%%

\subsection{Eksperyment nr n}
Eksperyment nr n zakładał, iż...\\
Dla dowolnej liczby $N \in \mathbb{N}$ funkcję 
$F_{N}:\mathbb{C}^{\:N}\!\rightarrow\mathbb{C}^{\:N}$
zdefiniowaną w następujący sposób:
\vspace{-0.4cm}
\begin{equation}
 \forall\:\mathbf{x} \in \mathbb{C}^{\:N}\:\:\:\:\:
 \forall\:k \in \{\:0,\dots,N - 1\:\!\}\:\:\:\:\:
 F_{N}\:\!(\:\mathbf{x}\:)_{k}\:\stackrel{\Delta}{=}\:
 \frac{1}{\sqrt{N}}\:
 \sum_{n = 0}^{N - 1}\:x_{n}\:\cdot\:
 e^{-j 2 \pi n k / N}
 \label{równanie dyskretnej transformaty Fouriera}
\end{equation}
nazywamy $N$ -- punktowym prostym jednowymiarowym dyskretnym przekształceniem Fouriera.
Na Rys. \ref{FFT} przedstawiono szybki algorytm obliczania dyskretnego 
przekształcenia Fouriera\footnote{ang. \textit{Fast Fourier Transform}}.

\subsubsection{Założenia}

\subsubsection{Przebieg}

\subsubsection{Rezultat}
\newpage

	\section{Wnioski}
	
	Spośród zaimplementowanych metod rekonstrukcji najlepsze wyniki uzyskano dla interpolacji sinc. Metoda ta osiąga najniższe wartości błędu średniokwadratowego (MSE) oraz najwyższe wartości współczynników SNR i PSNR.
	
	Metoda interpolacji liniowej (FOH) daje wyniki pośrednie – znacząco lepsze niż zero-order hold (ZOH), jednak gorsze niż rekonstrukcja sinc. W szczególności FOH redukuje efekt „schodków” charakterystyczny dla ZOH, ale nadal wprowadza zauważalne zniekształcenia.
	
	Metoda ZOH cechuje się największym błędem rekonstrukcji, co wynika z utrzymywania stałej wartości sygnału pomiędzy próbkami.
	
	W przypadku maksymalnej różnicy (MD) interpolacja sinc nie zawsze daje najlepsze wyniki względem FOH, co może wynikać z ograniczenia liczby próbek w implementacji.
	
	Zjawisko aliasingu potwierdza, że niespełnienie warunku Nyquista opisanego wzorem $f_s < 2\cdot f_{max}$ prowadzi do nieodwracalnej utraty informacji o sygnale.
	
	Przeprowadzone eksperymenty pokazują, że aliasing nie zależy od amplitudy sygnału, lecz wyłącznie od relacji między częstotliwością sygnału a częstotliwością próbkowania. Zmiana amplitudy wpływa jedynie na widoczność zjawiska na wykresach, ale nie eliminuje jego przyczyny.

	\renewcommand\refname{Bibliografia}
	\bibliographystyle{plain}
	\bibliography{bibliografia}

\end{document}
